\chapter{Generics}\label{ch-generics}

A stack of numbers and a stack of names are the same code. A function
that gives the first item of a list does the same whatever the items
are. \emph{Generics} let such code be written once, with the type left
open, and still be checked for every type it is used with.

\section{A generic function}

A function names the type it leaves open in square brackets after its
name, and uses that name where a type would go:

\begin{zeno}
function first[T](items: list[T]) -> T {
    return items[0]
}

print(first([3, 1, 2]), first(["ada", "alan"]))
let n = first([10, 20])
print(n + 1)
\end{zeno}

\verb|T| is a \emph{type parameter}. Read the heading as "first, for
any type T, takes a list of T and gives a T". At each call Zeno works
out what \verb|T| is, so \verb|first([10, 20])| is known to be an
\verb|int|, and \verb|n + 1| is checked as an int sum.

Without generics, \verb|first| would be written once for each type, or
once with no types at all, and then nothing that used its answer could
be checked.

\begin{why}
Zeno already writes "a type applied to a type" as \verb|list[int]|.
\verb|first[T]| is the same sentence, so there is one rule and no new
punctuation. The angle brackets of C++, Java and Rust would make
\verb|a < b > c| a puzzle for the parser and the reader; Go chose
square brackets for that reason. (Decision 3.16.)
\end{why}

A type parameter starts with a capital letter, so a reader can tell it
is a type without looking back at the heading:

\begin{zeno}[fails=yes]
function first[t](items: list[t]) -> t {
    return items[0]
}
\end{zeno}

\section{Several type parameters}

A function can leave several types open. Here \verb|convert| takes a
list of one type and a function that turns it into another:

\begin{zeno}
function convert[T, U](items: list[T], change: function(T) -> U) -> list[U] {
    var out: list[U] = []
    for item in items {
        out.append(change(item))
    }
    return out
}

let labels = convert([1, 2, 3], function(n) { return f"#{n}" })
print(labels, labels[0].upper())
\end{zeno}

Zeno works out \verb|T| from the list, then calls the function in its
head with an \verb|int| to find what it gives back, and that is
\verb|U|.

\section{A generic class}

A class can leave types open the same way. \verb|Stack[T]| is a stack
of anything, and \verb|Stack[str]| is a stack of text:

\begin{zeno}
class Stack[T] {
    items: list[T] = []

    function push(self, item: T) {
        self.items.append(item)
    }

    function pop(self) -> T {
        return self.items.pop()
    }

    function is_empty(self) -> bool {
        return len(self.items) == 0
    }
}

let names: Stack[str] = Stack()
names.push("ada")
names.push("grace")
print(names.pop().upper(), names.is_empty(), names)
\end{zeno}

\begin{zeno}
class Pair[A, B] {
    left: A
    right: B
}

let p = Pair(36, "Ada")
print(p, p.left + 1, p.right.upper())
\end{zeno}

Inside \verb|Pair[A, B]|, \verb|A| is the type parameter, even though
\verb|A| is also the unit of current: the name declared right there in
the heading is the nearest, and the nearest name wins. A quantity
inside it is still written \verb|3 A|.

\section{Where the types come from}

A type argument is never written at the call. There is no
\verb|first[int](xs)|. Zeno finds each type parameter in this order:

\begin{enumerate}
\item from the arguments;
\item from an argument that is a function, by calling it in its head
  with what the others gave;
\item from the place the result goes, such as a variable with a
  written type.
\end{enumerate}

The third is what a stack that starts empty needs: nothing in
\verb|Stack()| says what it will hold, so the variable says it, as
\verb|let names: list[str] = []| does for a list. If nothing says,
strict mode asks for it, and shows how:

\begin{zeno}[fails=yes, cmd=zeno check --strict {file}]
class Box[T] {
    items: list[T]
}
let box = Box([])
\end{zeno}

A written number waits for the other arguments to decide. With
\verb|level| a \verb|u8|, \verb|same(level, 9)| is a \verb|u8|, and
\verb|same(level, 300)| is an error, so a generic used for bytes never
needs \verb|u8(0)| where \verb|0| is meant. (Decision 3.16c.)

\section{The body is checked for each type}

\verb|largest| below uses \verb|>|, and nothing in its heading says
that \verb|T| must have one:

\begin{zeno}
function largest[T](items: list[T]) -> T {
    var best = items[0]
    for item in items {
        if item > best {
            best = item
        }
    }
    return best
}
print(largest([3, 9, 2]), largest(["pear", "fig", "apple"]))
\end{zeno}

Zeno checks the body once for each type the program uses it with. For
ints and for text it passes. For a class, which has no \verb|>|, the
mistake is reported inside \verb|largest|, with a note that names the
type and the line that asked for it:

\begin{zeno}[fails=yes]
class Point {
    x: int
}

function largest[T](items: list[T]) -> T {
    var best = items[0]
    for item in items {
        if item > best {
            best = item
        }
    }
    return best
}
print(largest([Point(1), Point(2)]))
\end{zeno}

\begin{why}
Checking each use is what C++ did for thirty years, and it finds every
mistake before the program runs, which is the promise that matters.
The note under the error is there so that the message points back to
the call. A generic may still say what it asks of \verb|T|, which
puts the message at the call instead: that is the next section.
(Decision 3.16d.)
\end{why}

\section{Saying what a type must be}

A type parameter may say what it asks of its type, with \verb|is| and
an interface, as a class says what it is. This is called a
\emph{bound}:

\begin{zeno}[file=largest.zeno, dir=bounds]
function largest[T is Ordered](items: list[T]) -> T {
    var best = items[0]
    for item in items {
        if item > best {
            best = item
        }
    }
    return best
}
print(largest([3, 9, 2]), largest(["pear", "apple"]), largest([3 V, 12 V]))
\end{zeno}

\verb|Ordered| comes with Zeno. A type is \verb|Ordered| when
\verb|<| works on two of it: every kind of number, a quantity, and
text. Now a list of points is refused at the call, which names the
type and what was asked of it, and the body is not even looked at:

\begin{zeno}[fails=yes]
class Point {
    x: int
}

function largest[T is Ordered](items: list[T]) -> T {
    var best = items[0]
    for item in items {
        if item > best {
            best = item
        }
    }
    return best
}
print(largest([Point(1), Point(2)]))
\end{zeno}

A bound can also be an interface of the program's own
(\hyperref[ch-interfaces]{Interfaces}). Then \verb|T| is any class
that says it is that interface:

\begin{zeno}[file=shapes.zeno, dir=bounds]
interface Shape {
    function area(self) -> float
}

class Circle is Shape {
    radius: float
    function area(self) -> float {
        return 3.14159 * self.radius * self.radius
    }
}

class Square is Shape {
    side: float
    function area(self) -> float {
        return self.side * self.side
    }
}

function total[S is Shape](shapes: list[S]) -> float {
    var sum = 0.0
    for s in shapes {
        sum += s.area()
    }
    return sum
}

let circles = [Circle(1.0), Circle(2.0)]
print(f"{total(circles):.2f}")
print(total([Square(2.0), Square(3.0)]))
\end{zeno}

\verb|total| takes a \verb|list[Circle]| as it is. A function that
asked for \verb|shapes: list[Shape]| would not: a list of circles is
not a list of shapes, since a square could then be put in it. The
bound says ``a list of one kind of shape'', and each kind gets its own
copy of \verb|total|.

A class and an enum take bounds too, as in
\verb|class Range[T is Ordered]|. Each type parameter has at most one
bound, and a comma always begins the next parameter, so
\verb|[T is Ordered, Shape]| would make \verb|Shape| a second type
parameter; it is refused with a message that says so. A bound is a
promise checked at the call, not all that the body may do: the body is
still checked for each copy, and a generic with no bound that uses
\verb|>| is still taken, as before. A generic that nothing calls is
checked with each bounded parameter standing for its interface, so a
call of a function the interface does not list is found even then. A
bound makes no code: the C of a program with bounds is the C of the
same program without them.

\begin{why}
Checking each copy finds every mistake, but it finds it inside the
generic, in words about \verb|>| that the caller did not write. A bound
moves the message to the call and puts the rule in the heading, where
a reader looks first. It is written with \verb|is| because a class
already says what it is that way, and an interface already names a set
of functions, so nothing new had to be learnt. A class cannot be
\verb|Ordered| yet: that waits for a class to say what \verb|<| means
for its objects. (Decisions 58.3 to 58.5.)
\end{why}

\section{A generic enum}

An enum can leave a type open too. The two most used are a value that
may be missing and a result that may be an error:

\begin{zeno}[file=maybe.zeno, dir=enums]
enum Maybe[T] {
    Nothing
    Just(value: T)
}

enum Result[T, E] {
    Ok(value: T)
    Failed(error: E)
}

function first_even(numbers: list[int]) -> Maybe[int] {
    for n in numbers {
        if n % 2 == 0 {
            return Just(n)
        }
    }
    return Nothing
}

function parse_volts(text: str) -> Result[float, str] {
    if text.ends_with("V") {
        return Ok(float(text.slice(0, len(text) - 1)))
    }
    return Failed(f"no unit in {text}")
}

for numbers in [[1, 3, 4], [5, 7]] {
    match first_even(numbers) {
        case Just(n) { print(f"found {n}") }
        case Nothing { print("none of them") }
    }
}
for text in ["3.3V", "12"] {
    match parse_volts(text) {
        case Ok(volts) { print(volts * 2.0) }
        case Failed(error) { print("error:", error) }
    }
}
\end{zeno}

The rules are those of a generic class. \verb|Maybe[int]| is a type,
a copy of the enum with \verb|T| standing for \verb|int|, and a
compiled program has one enum for each copy. Its cases are written by
their names alone, as the cases of any enum are, and their types are
worked out where they are used. A case that carries a value says what
\verb|T| is by what it carries: \verb|Just(3)| is a
\verb|Maybe[int]|. A case that carries nothing says nothing, so
\verb|Nothing| takes its type from the place it goes: the type of a
variable, a parameter, a field or a list, what the function returns,
or the other side of \verb|==|. Where nothing says, it is an error
that shows what to write:

\begin{zeno}[fails=yes]
enum Maybe[T] {
    Nothing
    Just(value: T)
}

let missing = Nothing
\end{zeno}

\verb|Ok(3)| says what \verb|T| is but not \verb|E|, so where the
result has no type of its own yet it is written
\verb|let r: Result[int, str] = Ok(3)|. A \verb|match| must handle
every case, as for any enum. A generic enum cannot give its cases
numbers, since a case with a number carries nothing.

\begin{why}
\verb|T or none| already says ``a \verb|T| or nothing'', and stays the
short way to say it. A generic enum is for the rest: a reason with the
nothing, two kinds of answer, or an enum that holds itself, such as
\verb|enum Chain[T] { End; Link(item: T, rest: Chain[T]) }|. It
needed no new rule, only the rules of a generic class applied to an
enum. Writing \verb|Nothing| alone and taking its type from where it
goes is what lets a case be written as any other enum's case is.
(Decisions 58.1 and 58.2.)
\end{why}

\section{A generic interface}

An interface can have type parameters, and a class says which copy of
it the class is:

\begin{zeno}[file=sources.zeno, dir=interfaces]
interface Source[T] {
    function next(self) -> T or none
}

class Countdown is Source[int] {
    left: int
    function next(self) -> int or none {
        if self.left == 0 {
            return none
        }
        self.left -= 1
        return self.left + 1
    }
}

class Repeat[T] is Source[T] {
    value: T
    times: int
    function next(self) -> T or none {
        if self.times == 0 {
            return none
        }
        self.times -= 1
        return self.value
    }
}

function drain(source: Source[int]) -> int {
    var sum = 0
    var item = source.next()
    while item is not none {
        sum += item
        item = source.next()
    }
    return sum
}

print(drain(Countdown(4)), drain(Repeat(5, 3)))
\end{zeno}

\verb|Source[int]| is a copy of the interface, as \verb|Stack[int]| is
a copy of a class, and every rule of an interface holds for it.
\verb|Countdown| is a \verb|Source[int]|; \verb|Repeat[int]| is one
too, and \verb|Repeat[str]| is a \verb|Source[str]|. The two copies
are two interfaces, so a \verb|Source[str]| does not go where a
\verb|Source[int]| is wanted. A copy can be a bound,
\verb|[S is Source[str]]|, and an \verb|owned interface Holder[T]|
has owned copies. A compiled program has one \verb|switch| for each
copy it uses, over the classes that are that copy. (Decision 58.6.)

\section{What a compiled program does}

In a compiled program each use becomes its own copy:
\verb|largest| for ints and \verb|largest| for text are two functions
in the C, each as fast as if it had been written by hand. Nothing is
looked up while the program runs. The cost is size: each type a
generic is used with adds a copy to the program, which on a
microcontroller is flash. That is why an untyped function is not made
generic by itself: a program asks for copies by writing \verb|[T]|.

\section{What generics leave out}

These are refused, each with a message that says so: a method with
type parameters of its own, a generic function or a case of a generic
enum kept as a value or passed along, the types written at the call
(there is no \verb|first[int](xs)|), and raising an object of a
generic class.

\begin{zeno}[fails=yes]
function first[T](items: list[T]) -> T {
    return items[0]
}
let pick = first
\end{zeno}

\verb|pick| would have to be every copy of \verb|first| at once. A
function written where a value goes can stand for one copy:
\verb|function(xs: list[int]) -> int { return first(xs) }|.
(Decisions 3.16e and 58.1.)

A method of \verb|class Box[T]| can use \verb|T|, but cannot have a
\verb|U| of its own. The message says what to write instead: a
function of the file that takes the object first, which does the same
work and which every build takes:

\begin{zeno}[file=boxes.zeno, dir=leave-out]
class Box[T] {
    item: T
}

function convert[T, U](box: Box[T], change: function(T) -> U) -> Box[U] {
    return Box(change(box.item))
}

let count = Box(42)
let text = convert(count, function(n: int) -> str { return f"#{n}" })
print(text.item)
\end{zeno}

(Decision 58.7.) And a build cannot tell what \verb|T| is from a list
whose items are not all one type, such as \verb|first([200 ms, 1 s])|,
since a compiled list has one type; \verb|zeno run| takes it. The
message says to write them all in one unit, such as \verb|ms|.
(Decision 58.8.)

\begin{summary}
\verb|function name[T](...)| and \verb|class Name[T]| leave a type
open. Type parameters start with a capital, and are worked out at each
call from the arguments, from a function argument, or from where the
result goes, never written at the call. The body is checked for every
type it is used with, and a compiled program has a copy for each.
\verb|[T is Ordered]| or \verb|[S is Shape]| says what a type must be,
and is checked at the call. \verb|enum Maybe[T]| and
\verb|interface Source[T]| leave a type open as a class does; a case
that carries nothing takes its type from where it goes, and a class
says which copy of an interface it is.
\end{summary}

\begin{tryit}
\begin{enumerate}
\item Write \verb|function last[T](items: list[T]) -> T or none|,
  which gives \verb|none| for an empty list.
\item Write a generic \verb|class Ring[T]| that holds at most
  \verb|size| items and drops the oldest when a new one comes. Use it
  for a ring of \verb|u8| and a ring of \verb|V|.
\item Write a generic \verb|count_if[T]| that takes
  \verb|items: list[T]| and \verb|test: function(T) -> bool|, and
  gives how many items pass the test.
\item Give \verb|largest| a list of \verb|bool| and read the message.
  Then take \verb|is Ordered| off its heading and read the message
  again. Where does each one point?
\item Write \verb|function or_else[T](m: Maybe[T], fallback: T) -> T|
  and call it as \verb|or_else(Nothing, 5)|. What tells \verb|Nothing|
  which \verb|Maybe| it is?
\item Write a class \verb|Letters is Source[str]| that gives the
  letters of a text one by one, and pass it to \verb|drain|. Read the
  message, then write a \verb|drain_text| for it.
\end{enumerate}
\end{tryit}
